\documentclass[11pt,a4paper]{article}
\usepackage[margin=2.2cm]{geometry}
\usepackage{fontspec}
\setmainfont{Latin Modern Roman}
\setmonofont{Latin Modern Mono}
\usepackage[provide=*]{babel}
\babelprovide[import,main]{romanian}
\usepackage{xcolor}
\usepackage{booktabs,tabularx,array}
\usepackage{graphicx}
\usepackage{enumitem}
\usepackage{listings}
\usepackage[most]{tcolorbox}
\usepackage[hidelinks]{hyperref}
\usepackage{titlesec}
\usepackage{fancyhdr}
\usepackage{needspace}

\definecolor{accent}{HTML}{1F5F8B}
\definecolor{trap}{HTML}{B3261E}

\titleformat{\section}{\Large\bfseries\color{accent}}{}{0pt}{}[\vspace{0.2em}\color{accent}\titlerule]
\titlespacing*{\section}{0pt}{1.6em}{0.8em}
\titleformat{\subsection}{\large\bfseries\color{accent}}{}{0pt}{}
\titlespacing*{\subsection}{0pt}{1.1em}{0.4em}
\setlist{itemsep=0.25em,topsep=0.3em}
\setlength{\parindent}{0pt}
\setlength{\parskip}{0.4em}

\pagestyle{fancy}
\fancyhf{}
\renewcommand{\headrulewidth}{0.3pt}
\fancyhead[L]{Administrarea rețelelor}
\fancyhead[R]{\cursnr}
\fancyfoot[C]{\thepage}

\lstset{basicstyle=\ttfamily\small,backgroundcolor=\color{black!4},frame=single,
  rulecolor=\color{black!20},xleftmargin=0.5em,xrightmargin=0.5em,
  columns=fullflexible,keepspaces=true,showstringspaces=false,breaklines=true}

\newtcolorbox{capcana}[1][Capcane de examen]{enhanced,colback=trap!5,colframe=trap,
  fonttitle=\bfseries,title={\,#1},boxrule=0.6pt,arc=2pt,left=6pt,right=6pt}

\newcommand{\img}[3][0.85]{\IfFileExists{#2}{\begin{center}
  \includegraphics[width=#1\linewidth]{#2}\\[-0.2em]{\small\itshape #3}\end{center}}{}}
\newcommand{\cod}[1]{\texttt{#1}}
\newcolumntype{Y}{>{\raggedright\arraybackslash}X}
\newcommand{\titlu}[2]{\begin{center}{\Huge\bfseries\color{accent} #1}\\[0.4em]{\large #2}\end{center}}

% titlurile nu rămân singure la capăt de pagină
\let\oldsection\section
\renewcommand{\section}{\Needspace{9\baselineskip}\oldsection}
\let\oldsubsection\subsection
\renewcommand{\subsection}{\Needspace{6\baselineskip}\oldsubsection}

\newcommand{\cursnr}{Curs 3 -- Hub, switch, ARP}
\begin{document}
\titlu{Hub, switch și ARP}{Administrarea rețelelor \textbullet{} Curs 3}

\section{Hub, switch (slide 45--48)}

\subsection*{Rețeaua veche cu cablu coaxial (în serie)}
\begin{itemize}
  \item \textbf{Single point of failure:} o problemă pe cablu sau la o placă de rețea strică comunicația pe toată rețeaua.
  \item Practic, PC-urile erau legate în serie.
  \item \textbf{Doar broadcast, necontrolat:} chiar dacă vorbesc doar 2 calculatoare, ceilalți interceptează.
  \item \textbf{Securitate:} un „spion” se injectează ușor pe cablu și vede tot. \textbf{Spoofing} ușor (impersonarea nodurilor).
  \item \textbf{Instabilă:} afectată de perturbații electromagnetice.
\end{itemize}

\subsection*{Rețea cu hub}
\begin{itemize}
  \item Coaxialul se înlocuiește cu \textbf{cablu cu 4 perechi torsadate $+$ conectori RJ-45} și un \textbf{hub}. Rezultă o rețea \textbf{stea}.
  \item \textbf{Rezolvă:} single point of failure (afectat e doar calculatorul cu probleme) și perturbațiile electromagnetice.
  \item \textbf{Hub $=$ echipament activ la L1 OSI}, doar nivel electric, fără „inteligență”.
  \item \textbf{NU rezolvă:} domeniul de broadcast (pachetul ajunge la toate calculatoarele, fie unicast, multicast sau broadcast) și spionarea (un cablu adițional în hub vede toată rețeaua, cu posibilitate de spoofing).
  \item Practic, dacă te conectai la oricare port, puteai să asculți ce vorbesc celelalte porturi.
\end{itemize}

\subsection*{Rețea cu switch}
\begin{itemize}
  \item \textbf{Switch $=$ echipament activ la L2 OSI.}
  \item Identifică \textbf{MAC-urile} de pe porturi și ține o \textbf{tabelă MAC} (temporară).
  \item Citește frame-ul Ethernet (MAC sursă și destinație) și trimite pachetul \textbf{doar pe portul} unde e MAC-ul destinație.
  \item Reduce broadcast-ul, îmbunătățește rata de transfer, \textbf{rezolvă spionarea unicast} și \textbf{parțial spoofing-ul}.
  \item Performanța depinde de \textbf{câte pachete Ethernet procesează pe unitatea de timp}, nu doar de viteza porturilor.
  \item Există switch-uri care își trimit între ele tabelele de rutare prin anumite protocoale.
  \item Contează mărimea tabelei și puterea cu care poate trimite switch-ul (câte pachete/sec).
  \item Avem:
  \begin{itemize}
    \item \textbf{Unmanaged:} doar funcția de bază, nu se configurează și nu se monitorizează, cost mic, rețele simple.
    \item \textbf{Managed:} configurabil și monitorizabil prin \textbf{serial UART, web, telnet, ssh}. Oferă facilități de securitate și stabilitate. Cost mai mare (enterprise mult mai mare).
  \end{itemize}
  \item Te poți conecta remote (telnet, SSH, Web).
  \item Când se aprinde becul pe portul respectiv, înseamnă că a reușit să comunice cu celălalt device și are adresa MAC.
\end{itemize}

\begin{capcana}
\begin{itemize}
  \item \textbf{Hub $=$ L1, switch $=$ L2.}
  \item Hub-ul rezolvă single point of failure, dar \textbf{nu} broadcast-ul și \textbf{nu} spionarea.
  \item Switch-ul rezolvă spionarea, dar spoofing-ul doar \textbf{parțial}.
  \item \textbf{Managed} înseamnă configurare (serial, web, telnet, ssh). \textbf{Unmanaged} nu se configurează.
\end{itemize}
\end{capcana}

\section{ARP (slide 50--54)}

\subsection*{Idee}
\begin{itemize}
  \item \textbf{ARP} descoperă \textbf{adrese MAC} pe baza IP-ului. C1 vrea să trimită un pachet IP către C2, dar pachetele IP merg în frame-uri Ethernet adresate cu MAC, deci C1 trebuie să știe MAC-ul lui C2.
  \item Paradigma \textbf{request-response}. Slide-ul îl plasează la \textbf{nivelul 3} și zice că e transportat de pachete \textbf{Ethernet de broadcast}. Pachetul are structură fixă.
\end{itemize}

\medskip
\renewcommand{\arraystretch}{1.35}
\small
\begin{tabularx}{\linewidth}{@{}p{3.1cm} Y Y@{}}
\toprule
 & \textbf{ARP Request} & \textbf{ARP Response}\\
\midrule
\textbf{OpCode} & 1 & 2\\
\textbf{Ethernet} & \textbf{broadcast}, MAC dest \cod{FF:FF:FF:FF:FF:FF} & \textbf{unicast} către cel care a întrebat\\
\textbf{Source MAC / IP} & cel care trimite cererea & cel care răspunde (IP $=$ Destination Protocol Address din request)\\
\textbf{Destination MAC} (în ARP) & \cod{00:00:00:00:00:00} (necunoscută) & MAC-ul celui care a trimis request-ul\\
\textbf{Destination IP} & IP-ul căruia i se cere MAC-ul & IP-ul celui care a trimis request-ul\\
\bottomrule
\end{tabularx}
\normalsize

\subsection*{Variante și mecanisme}
\begin{itemize}
  \item \textbf{ARP Probe:} request broadcast \textbf{înainte} de setarea IP-ului, ca să verifice dacă IP-ul e deja folosit în rețea. Dacă apare răspuns, IP-ul e deja ocupat.
  \item \textbf{GARP (Gratuitous ARP / ARP Announcement):} dispozitivul își anunță \textbf{neîntrebat} MAC-ul prin broadcast ARP Request. Se recunoaște prin \textbf{IP sursă $=$ IP destinație}.
  \item \textbf{ARP Proxy:} routerul răspunde la ARP Request pentru IP-uri din \textbf{altă rețea}, cu \textbf{propriul MAC}, și ține o tabelă internă pentru IP-urile din afară. Practic, router-ul răspunde cu MAC-ul lui pentru toate PC-urile sale.
  \item \textbf{ARP Caching:} dispozitivul ține o tabelă MAC-IP, ca să reducă cererile. O intrare se invalidează prin \textbf{timeout} sau când se constată că \textbf{nu mai merg conexiunile} cu perechea IP-MAC. Tabela se actualizează și la \textbf{GARP}.
  \item \textbf{ARP Spoofing:} un alt dispozitiv răspunde \textbf{mai repede} decât cel real, ca să îl înlocuiască în rețea.
\end{itemize}

\subsection*{Comenzi Linux}
\begin{itemize}
  \item \cod{arp -an}: afișează tabela. \cod{<incomplete>} în loc de MAC $=$ s-a cerut comunicația cu acel IP, dar nu s-a descoperit MAC-ul.
  \item \cod{arp -i eno1 -s 192.168.4.123 1C:00:34:22:11:AA}: adaugă manual o pereche IP-MAC.
  \item \cod{arp -d 192.168.4.123}: șterge o înregistrare.
  \item \cod{arping}: trimite forțat cerere ARP către un IP, așteaptă răspunsul și măsoară timpul. E ca \cod{ping}, dar \textbf{nu folosește ICMP}.
\end{itemize}

\begin{capcana}
\begin{itemize}
  \item \textbf{Request:} OpCode 1, Ethernet broadcast, dar Destination MAC din pachetul ARP $=$ \cod{00:..:00}.
  \item \textbf{Response:} OpCode 2, \textbf{unicast}.
  \item \textbf{GARP} se recunoaște după IP sursă $=$ IP destinație, nu după OpCode.
  \item \textbf{Probe} $=$ înainte de a seta IP-ul. \textbf{GARP} $=$ anunț după.
  \item \textbf{Proxy ARP} răspunde cu \textbf{MAC-ul routerului}.
  \item \textbf{\cod{arping} $\neq$ ICMP.}
  \item Intrarea ARP se invalidează prin \textbf{timeout} sau când conexiunea nu mai merge.
\end{itemize}
\end{capcana}
\end{document}
